\section{Discussion}\label{sec:discussion}
The results answer the paper's question in a way that depends on how realistic the substrate is,
and they do so consistently across all five disease types. Three threads run through them: how the
disease type shapes the epidemic, why targeting bridges beats targeting hubs, and why the choice of
substrate changes what containment is even possible.

\subsection{Disease type shapes the curve}
The five types separate cleanly by their dynamical structure rather than by any single rate. The
immunizing measles type produces the tall, right-skewed, single-peaked curve characteristic of
SIR: infectious cities find abundant susceptible neighbours, incidence grows, and then the
susceptible pool is depleted irreversibly, so the outbreak burns out~\cite{kermack:sir}. The
endemic gonorrhea type has no recovered compartment, so reinfection continually replenishes
susceptibles and incidence settles at an endemic equilibrium rather than declining to zero. The
three intermediate types interpolate as expected: latency (SEIR) delays and lowers the peak, and
both isolation (SEIQRD) and short infectious windows blunt it.

One coincidence deserves care. The waning-immunity influenza type and the isolation-controlled
Ebola type land at almost the same peak height ($\sim$$8$M on air), but for mechanistic rather than
biological reasons: influenza's basic reproduction number is low ($R_0\approx1.5$), while Ebola's
is higher ($R_0\approx2$) but its isolation shortens the effective infectious period
($1/(\gamma+\kappa)\approx3.3$~days), so the peak, which depends on the infectious period as well
as $R_0$, settles at a similar level. This tells us the model is internally coherent, not that an
Ebola outbreak would resemble a flu season. It would not: the isolation type is exactly where the
well-mixed metapopulation framing is weakest (Section~\ref{sec:limitations}), so its absolute
number is the least literal of the five, and it is run on the same footing as the others only for
the sake of comparison.

\subsection{Why bridges beat hubs}
Across every disease and every rung, immunising the highest-betweenness cities reduced the peak
more than immunising the highest-degree cities, and both beat random selection. The mechanism is
structural. Because transmission on a network follows paths, and betweenness counts how many
shortest paths run through a node, removing high-betweenness cities severs the routes that connect
otherwise weakly coupled regions~\cite{kitsak:spreaders}. Removing high-degree cities instead
deletes locally busy nodes whose traffic is often redundant, since their neighbours can still reach
each other by other short paths, so the giant component and its shortest-path structure survive
more intact. Betweenness targets the bottlenecks of reachability; degree targets a volume the
network can partly route around.

The realism ladder amplifies this. The degree--betweenness gap in the resulting peak roughly
doubles from the air-only network ($4$--$7.5\%$) to the full multilayer ($\approx$$12\%$) for every
disease. The added layers create few genuinely new bridges; instead they channel the extra
recurrent and ferry traffic through the same inter-community conduits, so a high-betweenness city
bridges several layers at once and removing it disrupts all of them~\cite{sorianopanos:multiplex}.
A single-layer model, which sees only air bridges, therefore \emph{understates} how much
bridge-targeting gains, the opposite of the intuition that redundant layers should dilute
targeting.

\subsection{Realism mostly adds people, and changes what containment means}
Climbing the ladder raised the absolute peaks, but the spread results
(Section~\ref{sec:res-spread}) show that most of that rise is a population effect: the land layer
brings inland, non-airport cities into the network, and the peak, as a fraction of the population,
stays nearly flat across rungs. This is why we report that fraction rather than the raw count,
which would let the land layer's extra population masquerade as extra severity. Ferries change
little, because they reach only a handful of coastal cities; land's real contribution is dynamical,
synchronising a larger connected population through recurrent commuting~\cite{balcan:recurrent}.

The interdiction results carry the sharpest practical implication. In an air-only model, grounding
all flights ends the outbreak; on the multimodal substrate the same closure leaves a third to a
half of the peak standing, and closing only the busiest hubs does almost nothing. Whether route
closure can contain an epidemic at all thus flips with the substrate. This echoes the broader
lesson that single-layer transport models systematically mislead about containment once ground and
maritime mobility are included~\cite{balcan:multiscale}, and it is the clearest argument in the
paper for building the multilayer substrate rather than the air network alone.

\subsection{Limitations}\label{sec:limitations}
The study is comparative, not predictive, and its absolute counts should not be read as forecasts.
Three limitations bound the interpretation. First, the well-mixed metapopulation makes transmission
proportional to the local infectious fraction, which fits casual-contact respiratory and endemic
diseases far better than a contact-limited pathogen such as Ebola: real filovirus transmission
requires direct contact with symptomatic or deceased cases and is shaped by care settings that a
constant isolation rate only crudely approximates. Its peak is therefore the least trustworthy of
the five, and our metric, peak active infection, does not even reflect the isolation type's high
case-fatality ($\mu\approx0.71$), so a modest infectious peak can still be devastating. Second,
each configuration is a single deterministic run at fixed coverage, efficacy, and seed; a
stochastic ensemble and a coverage sweep would put confidence bounds on the reductions, though the
strategy \emph{ordering} is stable enough that we do not expect it to change. Third, parameters are
justified from the literature rather than calibrated to any outbreak, and the transport weights are
proxies for passenger volume, so the results are internally consistent comparisons, not fits to
data.

These bounds are also the roadmap. The same engine runs stochastic ensembles, other regions, and
finer coverage grids without modification, and the interactive simulator turns each of those into a
configuration change rather than a code change, so the findings here, bridges over hubs, realism
amplifying that gap, and the limits of route closure, form a stable core on which a more
quantitative, region-specific study can build.
